% Lean4ProofTrees.tex — natural-deduction proof trees paired with Lean 4 terms.
% Compile:  scripts/tex2pdf.sh scripts/Lean4ProofTrees.tex  ->  generated-latex/Lean4ProofTrees.pdf
\documentclass[11pt]{article}
\usepackage[margin=1in]{geometry}
\usepackage{amsmath, amssymb}
\usepackage{bussproofs}
\usepackage{fontspec}
\setmonofont{Menlo}[Scale=0.85]

\title{Lean 4 proof trees}
\author{ScottLean4 \textperiodcentered\ for Dana Scott}
\date{}
\newcommand{\lean}[1]{\texttt{#1}}

\begin{document}
\maketitle

By the Curry--Howard correspondence, a natural-deduction \emph{proof tree} is the
same object as a Lean~4 \emph{proof term}: each inference rule is a term
constructor, and discharging an assumption is a $\lambda$-abstraction.

\bigskip
\noindent\textbf{$\wedge$-introduction.}\quad From $a$ and $b$, conclude $a\wedge b$.
\begin{center}
\AxiomC{$a$}
\AxiomC{$b$}
\RightLabel{\scriptsize $\wedge$I}
\BinaryInfC{$a \wedge b$}
\DisplayProof
\end{center}
Lean: \lean{⟨ha, hb⟩}\; (i.e.\ \lean{And.intro ha hb}).

\bigskip
\noindent\textbf{$\rightarrow$-elimination (modus ponens).}\quad
From $a \to b$ and $a$, conclude $b$.
\begin{center}
\AxiomC{$a \to b$}
\AxiomC{$a$}
\RightLabel{\scriptsize $\to$E}
\BinaryInfC{$b$}
\DisplayProof
\end{center}
Lean: \lean{h ha}\; (apply the arrow to its argument; the \lean{apply} tactic).

\bigskip
\noindent\textbf{$\rightarrow$-introduction / discharge.}\quad
Proof of $a \to (b \to a)$; superscripts mark the discharged assumptions.
\begin{center}
\AxiomC{$[a]^{1}$}
\RightLabel{\scriptsize $\to$I$^{2}$}
\UnaryInfC{$b \to a$}
\RightLabel{\scriptsize $\to$I$^{1}$}
\UnaryInfC{$a \to (b \to a)$}
\DisplayProof
\end{center}
Lean: \lean{fun (ha : a) (\_ : b) => ha}.

\bigskip
\noindent\textbf{A two-branch derivation.}\quad
$(a \wedge b) \to (b \wedge a)$.
\begin{center}
\AxiomC{$[a \wedge b]^{1}$}
\RightLabel{\scriptsize $\wedge$E$_2$}
\UnaryInfC{$b$}
\AxiomC{$[a \wedge b]^{1}$}
\RightLabel{\scriptsize $\wedge$E$_1$}
\UnaryInfC{$a$}
\RightLabel{\scriptsize $\wedge$I}
\BinaryInfC{$b \wedge a$}
\RightLabel{\scriptsize $\to$I$^{1}$}
\UnaryInfC{$(a \wedge b) \to (b \wedge a)$}
\DisplayProof
\end{center}
Lean: \lean{fun h => ⟨h.right, h.left⟩}.

\end{document}
